\section{Discussion and Limitations}
\label{sec:discussion}

The results presented in Section~\ref{sec:evaluation} demonstrate that \texttt{LLMOrch-VAPT} provides a robust and cost-effective foundation for autonomous security validation. However, the deployment of such systems in enterprise environments raises several operational, ethical, and technical questions that warrant further discussion.

\subsection{Operational Continuity and SOC Integration}
A primary advantage of the \texttt{LLMOrch-VAPT} architecture is its readiness for integration into existing Security Operations Center (SOC) workflows. By providing a provider-agnostic abstraction layer, organizations can maintain a "Red Team Agent" fleet that is resilient to the volatile landscape of AI service providers. Furthermore, the stateful checkpointing mechanism used for APF (Section~\ref{sec:apf}) enables long-running engagements that can span across shift changes and maintenance windows without losing progress, a critical requirement for industrial VAPT.

\subsection{Ethical Considerations and Dual-Use Mitigation}
As with any autonomous exploitation framework, \texttt{LLMOrch-VAPT} possesses significant dual-use potential. While its primary mission is to assist defenders in continuous security validation, the same orchestration logic could be utilized by malicious actors to scale automated attacks. 

To mitigate these risks, we have prioritized \textit{safety governance} as a core architectural component. Our 100\% adherence to HITL gates (Section~\ref{sec:evaluation}) demonstrates that autonomous agents can be effectively constrained. We advocate for a "Safety-First" deployment model where destructive capabilities are strictly authenticated and audited. Furthermore, by open-sourcing the platform, we aim to democratize access to high-end defensive automation, helping to balance the asymmetric advantage currently held by well-resourced adversarial groups.

\subsection{Technical Limitations}
Despite the performance gains, \texttt{LLMOrch-VAPT} is subject to several technical limitations:
\begin{itemize}
    \item \textbf{Context Window Bottlenecks}: While the Master-Worker hierarchy reduces the context load on individual agents, extremely large-scale network recons can still exceed the context windows of smaller "Flash" models, potentially leading to reasoning loss.
    \item \textbf{Business Logic Vulnerabilities}: Current LLM-based agents, including our specialized workers, still struggle with complex business logic vulnerabilities that require deep, domain-specific understanding of application state and user intent.
    \item \textbf{Adversarial Defenses}: The system's reliance on technical signals for DCAT (Section~\ref{sec:dcat}) makes it potentially susceptible to "signal spoofing" by sophisticated defenders who attempt to bait the orchestrator into using more expensive (and potentially slower) reasoning tiers.
\end{itemize}

\subsection{Future Work}
Future research directions include the integration of \textit{experience-based RAG} to allow agents to learn from successful exploit chains over time, and the development of \textit{adversarial prompt protection} to ensure the integrity of the orchestrator's decision-making graph against injection attacks.
